\documentclass[12pt, xcolor=dvipsnames,svgnames,x11names]{article}
\usepackage{xcolor}
\usepackage{xspace}
\usepackage{url}
\usepackage[colorlinks=true,bookmarks=false,linkcolor=black,urlcolor=blue,citecolor=black]{hyperref}
\usepackage{fancyhdr}
\usepackage[top=1in,bottom=1.2in,left=1in,right=1in]{geometry}
\usepackage{multicol}
\usepackage{pgfplots}
\usepackage{tikz}
\usetikzlibrary{circuits.logic.US, patterns}
\usepackage{mathpazo}
\usepackage{biolinum} 
\usepackage[all]{tcolorbox}
\usepackage{derivative}
\usepackage{../Lectures/rahul_math}
\renewcommand{\familydefault}{\sfdefault}

% --- Custom Color Palette from Image ---
\definecolor{headerRed}{HTML}{A3403D} % Dark red from image top
\definecolor{patternBg}{HTML}{F2D9D7} % Light pinkish bg from image bottom
\definecolor{binaryText}{HTML}{D1A7A4} % Muted color for binary digits

% --- Background Pattern Setup ---
\usepackage[pages=all]{background}
\backgroundsetup{
scale=1,
color=black,
opacity=1,
angle=0,
contents={%
  \begin{tikzpicture}[remember picture,overlay]
    % Top Header Bar
    \fill[headerRed] (current page.north west) rectangle ([yshift=-1.2cm]current page.north east);
    % Bottom Binary Background
    \fill[patternBg] (current page.south west) rectangle ([yshift=3.5cm]current page.south east);
    % Decorative Line (from image)
    \draw[headerRed, line width=2pt] ([xshift=1in, yshift=-2.5cm]current page.north west) -- ([xshift=8in, yshift=-2.5cm]current page.north west);
    % Binary Text Overlay (Sample)
    \node[anchor=south, binaryText, font=\ttfamily\scriptsize, inner sep=0pt, yshift=1.5cm] at (current page.south) {
        \begin{tabular}{c@{\hspace{0.5em}}c@{\hspace{0.5em}}c@{\hspace{0.5em}}c@{\hspace{0.5em}}c@{\hspace{0.5em}}c@{\hspace{0.5em}}c}
        0 0 0 0 1 1 1 1 0 0 0 0 1 1 0 1 1 1 0 0 0 1 \\
        1 0 0 1 0 1 0 0 0 1 0 0 1 1 0 0 1 0 1 1 \\
        0 0 1 1 0 0 1 0 1 1 1 0 1 1 0 1 0 1 0 0 0 \\
        The University of Alabama in Huntsville
        \end{tabular}
    };
  \end{tikzpicture}
}
}

\usetikzlibrary{positioning, arrows.meta, shadows}

% --- Header Styling ---
\makeatletter
\renewcommand{\maketitle}{%
    \vspace*{0.2cm}
    \begin{center}
        \begin{tcolorbox}[
            enhanced,
            colback=white,
            colframe=headerRed,
            arc=0mm,
            title={\textcolor{white}{\bfseries \@title}},
            coltitle=white,
            fonttitle=\bfseries\large,
            attach boxed title to top left={xshift=10pt, yshift=-\tcboxedtitleheight/2},
            boxed title style={sharp corners, colback=headerRed, frame hidden}
        ]
        \large
        \vspace{10pt}
        \noindent \textbf{Student Name:} \underline{\hspace{7cm}} \hfill \textbf{Points:} 30
        
        \vspace{10pt}
        \noindent \textbf{A Number:} \hfill \underline{\hspace{7cm}} \hfill \textbf{Due:} Feb 17, 2026
        \end{tcolorbox}
    \end{center}
    \vspace{10pt}
}
\makeatother

\newcommand{\hw}{Classwork 03//Computational Thinking, Convolutional Neural Net}
\newcommand{\duedate}{Jan 22, 2026}

\pagestyle{fancy}
\fancyhf{}
\lhead{\small CPE 487/587, Spring 2026}
\rhead{\small \textit{\hw}}
\cfoot{\thepage}
\renewcommand{\headrulewidth}{0pt}

\usepackage{booktabs} % For professional lines


% Define a very light gray for the rows
\definecolor{tablegray}{gray}{0.95}

\begin{document}

\title{\hw}
\maketitle




\section{Log (1+x) Layer \hfill \textbf{(10 Points)}}

Provide a code-snippet of how would you design a neural network Layer using PyTorch by inheriting from \texttt{nn.Module} that performs $\log(1+x)$ transformation on the feature matrix.

\subsection{Solution}

          \begin{minted}[
                linenos,
                frame=lines,
                framesep=2mm,
                bgcolor=pink!10,
                fontsize=\footnotesize,
                breaklines
            ]{Python}
class LogLayer(nn.Module):
    def __init__(self):
        super().__init__()
        
    def forward(self, x):
        logx = torch.log(x + 1)
        return logx
        
model = LogLayer()
\end{minted}

\section{Convolution Dimension\hfill \textbf{(5 Points)}}

Consider a kernel of dimension $3\times 3$ operating on an input image dimension of $5\times 5$ with a stride of 1. Consider no padding, what is the dimension of the output??


\subsection{Solution}

$3 \times 3$.

\section{Convolution Computation \hfill \textbf{(10 Points)}}

Given input matrix  as

\begin{align*}
I = \begin{bmatrix}
1 & 2 & 3 & 0 \\
     4 & 5 & 6 & 1\\
     7 & 8 & 9 & 2\\
     0 & 1 & 2 & 3
\end{bmatrix}
\end{align*}

and kernel as

\begin{align*}
K = \begin{bmatrix}
1 & 0 & -1\\
    1 & 0 & -1\\
     1 & 0 & -1
\end{bmatrix}
\end{align*}

Calculate the output feature map using convolution with a stride of 1.

\subsection{Solution}

\begin{align*}
Output = \begin{bmatrix}
-6 & 12\\
-6 & 8
\end{bmatrix}
\end{align*}



\section{Max Pooling \hfill \textbf{(5 Points)}}
Apply max pooling with a stride of 2 to 

\begin{align*}
I = \begin{bmatrix}
 3 & 2 & 5 & 1\\
     4 & 1 & 6 & 2\\
     1 & 7 & 3 & 4\\
     2 & 5 & 1 & 3
\end{bmatrix}
\end{align*}

\subsection{Solution}

\begin{align*}
O = \begin{bmatrix}
4 & 6\\
7 & 4
\end{bmatrix}
\end{align*}



\end{document}